\documentclass[sigconf,nonacm]{acmart}

\bibliographystyle{alpha}
\title{Persistence Diagrams as Barycenters for Wassertein Distances\\{\small Report on \textit{Wasserstein 
      Dictionaries of Persistence Diagrams}}}

\author{Matthieu Pierre Boyer}
\authornote{All authors contributed equally.}
\affiliation{%
	\institution{École Normale Supérieure}
	\city{Paris}
	\country{France}}
\email{matthieu.boyer@ens.fr}

\author{Gary Klajer}
\authornotemark[1]
\affiliation{%
	\institution{ÉNS Paris-Saclay}
	\city{Gif-sur-Yvette}
	\country{France}}
\email{gary.klajer@ens-paris-saclay.fr}

\author{Marc Perlade}
\authornotemark[1]
\affiliation{%
	\institution{École Normale Supérieure}
	\city{Paris}
	\country{France}}
\email{marc.perlade@ens.fr}

\def\abs#1{\left|#1\right|}


\begin{abstract}
  In this report based on \cite{sisouk2023wassersteindictionariespersistencediagrams},
  we study the idea of doing non-linear dimensionality reduction on a large dataset of persistence diagrams,
  using atomic persistence diagrams that span the dataset for the hull defined by the optimal transport of
  diagrams.
  We will first give the mathematical intuition on why this could work then provide an algorithm that
  proposes a method to compute such a representation.
	We will finally provide benchmarks on the different methods, implemented in Rust.
	% for a blazingly fast™ result.
\end{abstract}

\keywords{Persistence, Dictionaries, Optimization, Topological Data Analysis}

\begin{teaserfigure}
	% ce serait bien d'avoir une image ici
	%   \includegraphics[width=\textwidth]{sampleteaser}
	\caption{}
	\Description{}
	\label{fig:teaser}
\end{teaserfigure}

\begin{document}
\maketitle

\section*{Introduction}

\section{Mathematical preliminaries}
\subsection{Representation of persistence diagrams}
We are given a set $X = \left\{X_{1}, \ldots, X_{n}\right\}$ of $N$ persistence diagrams.
Each input diagram is given as the set of its persistence pairs mapped to points in the $2$-dimensional
plane\footnote{Or half-plane above the line $y = x$.}, along with their birth and death persistence.

\subsection{\texorpdfstring{$2$-dimensional}{2-dimensional} optimal transport}
Optimal transport is the mathematical field concerned with mapping similarities from a space with
some sense of geometry to the space of probability distributions on that space.


\subsection{}

\section{Algorithm}
\subsection{}

\subsection{Implementation}

\section{Analysis and comparison of implementations}

\appendix
\bibliography{report}
\end{document}
